\documentclass[12pt,a4paper]{article}

\usepackage[a4paper,margin=18mm]{geometry}
\usepackage[T2A]{fontenc}
\usepackage[utf8]{inputenc}
\usepackage[russian]{babel}
\usepackage{amsmath,amssymb,mathtools}
\usepackage{array,booktabs}
\usepackage{xcolor}
\usepackage{hyperref}
\usepackage{tikz}

\hypersetup{
  colorlinks=true,
  linkcolor=blue!55!black,
  urlcolor=blue!55!black
}

\setlength{\parindent}{0pt}
\setlength{\parskip}{0.55em}
\setlength{\tabcolsep}{3pt}
\renewcommand{\arraystretch}{1.25}

\begin{document}

% Краткое сообщение Володе:
% Володя, в четвёртом задании не дописан смысл третьего выражения: это сумма длин трёх сторон прямоугольника. В девятом задании при раскрытии скобок перед ними стоит минус, поэтому знак у семёрки тоже должен измениться; именно это повлияло на итоговое значение.

\begin{titlepage}
\centering
\vspace*{\fill}

{\LARGE\bfseries Ревью домашней работы\par}

\vspace{1.2cm}

{\large Автор: Лёвин Артём Александрович\par}

\vspace{0.6cm}

{\large 16 сентября 2026 года\par}

\vspace*{\fill}

\begin{tikzpicture}
\draw[line width=0.7pt] (0,0) .. controls (2,-0.25) and (4,0.25) .. (6,0);
\end{tikzpicture}

\end{titlepage}

\newpage
\section*{Сводная таблица}

{\small
\begin{center}
\begin{tabular}{p{0.09\linewidth} p{0.17\linewidth} p{0.25\linewidth} p{0.18\linewidth} p{0.18\linewidth}}
\toprule
\textbf{№ задания} & \textbf{Тема} & \textbf{Суть ошибки} & \textbf{Ваш ответ} & \textbf{Правильный ответ} \\
\midrule
\hyperref[task:4]{4} &
Геометрический смысл выражений &
Не объяснён смысл выражения $m+2n$ &
$m+2n$ --- не указано &
Сумма длин трёх сторон прямоугольника: одной стороны длины $m$ и двух сторон ширины $n$ \\
\addlinespace
\hyperref[task:9]{9} &
Раскрытие скобок и приведение подобных слагаемых &
При раскрытии $-(x+7)$ знак перед $7$ не изменён &
$18$ &
$4$ \\
\bottomrule
\end{tabular}
\end{center}
}

\newpage
\vspace*{\fill}
\begin{center}
{\LARGE\bfseries Разбор ошибок}
\end{center}
\vspace*{\fill}
\newpage

\phantomsection\label{task:4}
\section*{Задание 4}

\subsection*{Текст задания}
У прямоугольника длина равна $m$, ширина равна $n$. Объясните словами смысл каждого выражения:
\[
mn,\qquad 2m+2n,\qquad m+2n.
\]

\subsection*{Ваше решение}
Записано:
\[
mn\;\text{--- площадь},
\]
\[
2m+2n\;\text{--- периметр}.
\]
Для выражения
\[
m+2n
\]
смысл не указан: рядом поставлен вопросительный знак.

\subsection*{Ваш ответ}
Для первых двух выражений смысл указан, для $m+2n$ ответ не дан.

\subsection*{Ошибка}
\textcolor{red}{Ошибка:} не выполнена последняя часть задания --- не объяснён геометрический смысл выражения $m+2n$.

\subsection*{Где именно возникает ошибка}
Последний полностью правильный шаг:
\[
2m+2n\;\text{--- периметр прямоугольника}.
\]
Это верно, потому что у прямоугольника две стороны имеют длину $m$ и ещё две стороны имеют длину $n$, поэтому сумма длин всех четырёх сторон равна
\[
m+n+m+n=2m+2n.
\]

Первый неполный шаг:
\[
m+2n\;\text{--- ?}
\]
Здесь вычислять ничего не требуется: нужно связать каждое слагаемое с конкретными сторонами прямоугольника и словами описать получившуюся сумму длин.

\subsection*{Почему это неверно}
Задание просит объяснить смысл \textbf{каждого} из трёх выражений. Поэтому правильные объяснения для $mn$ и $2m+2n$ не делают решение полным, если третье выражение осталось без интерпретации.

У прямоугольника имеются две стороны длины $m$ и две стороны длины $n$. В выражении
\[
m+2n
\]
слагаемое $m$ обозначает длину одной из двух длинных сторон, а $2n=n+n$ --- сумму длин двух сторон ширины. Значит,
\[
m+2n=m+n+n.
\]
Это сумма длин трёх сторон прямоугольника: одной стороны длины $m$ и обеих сторон длины $n$. Иными словами, это длина трёх сторон прямоугольника, если одну из сторон длины $m$ не учитывать.

\textcolor{blue!60!black}{Ключевая идея:} чтобы объяснить геометрический смысл алгебраического выражения, нужно сопоставить каждое слагаемое реальному геометрическому объекту. Здесь $m$ и $n$ --- длины сторон, поэтому их суммы описывают суммы длин определённых сторон прямоугольника.

\subsection*{Исправление от последнего правильного шага}
После верной записи
\[
2m+2n\;\text{--- периметр}
\]
нужно рассмотреть оставшееся выражение:
\[
m+2n=m+n+n.
\]
Следовательно,
\[
\boxed{m+2n}
\]
--- сумма длин трёх сторон прямоугольника: одной стороны длины $m$ и двух сторон длины $n$.

\subsection*{Полное правильное решение}
Рассмотрим каждое выражение отдельно.

Первое выражение:
\[
mn.
\]
Площадь прямоугольника равна произведению его длины на ширину. Поэтому
\[
\boxed{mn\;\text{--- площадь прямоугольника}.}
\]

Второе выражение:
\[
2m+2n.
\]
У прямоугольника две стороны длины $m$ и две стороны длины $n$, поэтому сумма длин всех четырёх сторон равна
\[
m+m+n+n=2m+2n.
\]
Следовательно,
\[
\boxed{2m+2n\;\text{--- периметр прямоугольника}.}
\]

Третье выражение:
\[
m+2n=m+n+n.
\]
Оно содержит длину одной стороны $m$ и длины двух сторон $n$. Поэтому
\[
\boxed{m+2n}
\]
--- сумма длин трёх сторон прямоугольника, если не учитывать одну из сторон длины $m$.

\subsection*{Проверка}
Возьмём, например,
\[
m=5,\qquad n=3.
\]
Тогда
\[
mn=5\cdot3=15,
\]
что действительно является площадью прямоугольника со сторонами $5$ и $3$.

Далее,
\[
2m+2n=2\cdot5+2\cdot3=10+6=16,
\]
а сумма всех четырёх сторон равна
\[
5+3+5+3=16.
\]

Наконец,
\[
m+2n=5+2\cdot3=11,
\]
и сумма одной стороны длины $5$ и двух сторон длины $3$ равна
\[
5+3+3=11.
\]
Интерпретация согласуется с геометрией прямоугольника.

\subsection*{Итог}
\textcolor{teal!60!black}{Итог:} первые два объяснения были верными. Для завершения задания нужно было дописать, что $m+2n$ обозначает сумму длин трёх сторон прямоугольника: одной стороны длины $m$ и двух сторон длины $n$.

\subsection*{Задача на закрепление}
У прямоугольника длина равна $a$, ширина равна $b$. Объясните словами геометрический смысл выражения
\[
2a+b.
\]

\newpage
\phantomsection\label{task:9}
\section*{Задание 9}

\subsection*{Текст задания}
Упростите выражение и затем при $x=4$ найдите его значение:
\[
3(2x-5)-(x+7)+2(x-1).
\]

\subsection*{Ваше решение}
По записи решение имеет вид
\[
6x-15-x+7+2x-2=7x-10.
\]
Затем при $x=4$ получено
\[
7\cdot4-10=18.
\]

\subsection*{Ваш ответ}
\[
18.
\]

\subsection*{Ошибка}
\textcolor{red}{Ошибка:} при раскрытии скобок в выражении $-(x+7)$ знак перед числом $7$ не изменён.

\subsection*{Где именно возникает ошибка}
Правильно раскрыты первая и третья скобки:
\[
3(2x-5)=6x-15,
\]
\[
2(x-1)=2x-2.
\]

Ошибка появляется при переходе
\[
-(x+7)\longrightarrow -x+7.
\]
Это первый достоверно неверный шаг.

Правильный переход:
\[
-(x+7)=-x-7.
\]

\subsection*{Почему это неверно}
Минус перед скобками означает умножение \textbf{всего} выражения в скобках на $-1$:
\[
-(x+7)=(-1)(x+7).
\]
При умножении каждого слагаемого на $-1$ знак каждого слагаемого меняется:
\[
(-1)\cdot x=-x,
\qquad
(-1)\cdot7=-7.
\]
Поэтому после раскрытия скобок должно получиться
\[
-x-7,
\]
а не
\[
-x+7.
\]

Из-за этой ошибки постоянная часть выражения изменилась на $14$: вместо $-7$ было записано $+7$. Поэтому дальнейшее приведение подобных слагаемых и подстановка $x=4$ уже выполнялись для другого выражения и привели к неверному результату.

\textcolor{blue!60!black}{Ключевая идея:} если перед скобками стоит минус, при раскрытии скобок нужно изменить знак \textbf{у каждого} слагаемого внутри скобок.

\subsection*{Исправление от последнего правильного шага}
Сохраняем правильно раскрытые первую и третью скобки и исправляем только среднюю часть:
\[
3(2x-5)-(x+7)+2(x-1)
\]
\[
=6x-15-x-7+2x-2.
\]
Теперь приводим подобные слагаемые. Слагаемые с $x$ дают
\[
6x-x+2x=7x.
\]
Постоянные слагаемые дают
\[
-15-7-2=-24.
\]
Следовательно,
\[
\boxed{7x-24}.
\]

При $x=4$:
\[
7\cdot4-24=28-24=4.
\]
Значит,
\[
\boxed{4}.
\]

\subsection*{Полное правильное решение}
Дано выражение
\[
3(2x-5)-(x+7)+2(x-1).
\]

Раскроем первые скобки. Число $3$ умножается на каждое слагаемое внутри скобок:
\[
3(2x-5)=3\cdot2x-3\cdot5=6x-15.
\]

Перед вторыми скобками стоит минус. Это означает умножение каждого слагаемого в скобках на $-1$:
\[
-(x+7)=-x-7.
\]

Раскроем третьи скобки:
\[
2(x-1)=2x-2.
\]

Получаем
\[
6x-15-x-7+2x-2.
\]

Соберём подобные слагаемые с $x$:
\[
6x-x+2x=(6-1+2)x=7x.
\]

Соберём числа:
\[
-15-7-2=-24.
\]

Итак, упрощённое выражение равно
\[
\boxed{7x-24}.
\]

Теперь подставим $x=4$:
\[
7\cdot4-24=28-24=4.
\]

Следовательно, значение исходного выражения при $x=4$ равно
\[
\boxed{4}.
\]

\subsection*{Проверка}
Проверим результат прямой подстановкой $x=4$ в исходное выражение:
\[
3(2\cdot4-5)-(4+7)+2(4-1).
\]
Вычисляем значения в скобках:
\[
2\cdot4-5=8-5=3,
\]
\[
4+7=11,
\]
\[
4-1=3.
\]
Тогда
\[
3\cdot3-11+2\cdot3=9-11+6=4.
\]
Получили то же значение:
\[
4.
\]
Проверка подтверждает исправленное решение.

\subsection*{Итог}
\textcolor{teal!60!black}{Итог:} ошибка возникла при раскрытии скобок со знаком минус. Нужно было заменить $+(7)$ внутри скобок на $-7$ после раскрытия. Правильно упрощённое выражение равно $7x-24$, а при $x=4$ его значение равно $4$.

\subsection*{Задача на закрепление}
Упростите выражение и затем при $x=3$ найдите его значение:
\[
2(3x-4)-(x+5)+3(x-2).
\]

\vfill
\begin{center}
\small Лёвин Артём Александрович эксклюзивно для Володи
\end{center}

\end{document}
